\subsection{Discussion of Potential Data Leakage}
\label{sec:discussion_utbench}

The issue of data leakage, as discussed by Zhang et al.~\cite{zhang2023critical}, poses a critical risk to the evaluation of LLMs in the SE community.
In this section, we concentrate on the test generation task, for which we employ Defects4J as an evaluation dataset. 
To mitigate potential data leakage risk, we test LLMs' performance on a real-world industrial dataset\footnote{Due to the confidential policy of the company, we hide the information of internal subjects.}.
Unlike Defects4J, which partially relies on JUnit 3, the internal dataset adopts a unified Maven architecture with the JUnit 4 testing framework. This setup aligns with mainstream development practices and provides a more accurate measure of LLMs' capability to generate test cases for contemporary Java projects.



\begin{table*}[htbp]
\centering
\small
\caption{Performance of different LLMs on the internal dataset. The best results among all LLMs are \textbf{bolded}, and the second-best results are \uline{underlined}.}
\label{tab:discussion_utbench}
\begin{tabular}{lccccc}
\toprule
\textbf{LLMs} & \textbf{Correct} & \textbf{Passing} & \textbf{Failing} & \textbf{Build Error} & \textbf{Syntax Error} \\
\midrule
UniXcoder         & 5.39\%  & 6.59\%  & \textbf{1.95\%}  & 74.10\%  & 17.37\% \\
CodeT5-large      & 14.07\% & 16.32\% & 10.03\% & 66.77\%  & 6.89\%  \\
CodeT5p-770m      & 18.56\% & 19.61\% & 5.54\%  & 71.11\%  & \uline{3.74\%}  \\
DeepSeek-Coder-6b & 31.14\% & 31.59\% & 2.84\%  & \textbf{29.64\%}  & 35.93\% \\
CodeLlama-7b      & 25.00\% & 25.45\% & \uline{2.10\%}  & 34.28\%  & 38.17\% \\  
StarCoder2-7b     & 23.05\% & 24.40\% & 6.59\%  & 42.66\%  & 26.35\% \\
StarCoderBase-7b  & 23.65\% & 27.10\% & 2.40\%  & 42.96\%  & 27.54\% \\
CodeGen-6b        & 23.50\% & 23.95\% & \uline{2.10\%}  & 35.78\%  & 38.17\% \\
\midrule
GPT-3.5           & \textbf{52.54\%} & \textbf{54.34\%} & 12.43\% & \uline{32.04\%}  & \textbf{1.20\%}  \\
Llama3.1-8b       & 21.71\% & 22.01\% & 17.96\% & 51.50\%  & 8.53\% \\
Llama3.1-70b      & \uline{51.35\%} & \uline{52.40\%} & 10.48\% & {32.63\%}  & 4.49\% \\
Llama3-8b         & 28.44\% & 29.94\% & 10.33\% & 54.34\%  & 5.39\%  \\
Qwen2-7b          & 32.34\% & 32.49\% & 10.03\% & {50.90\%}  & 6.59\% \\
Qwen2-72b         & 42.37\% & 44.61\% & 7.63\% & 42.37\%  & 5.39\%  \\
\bottomrule
\end{tabular}
\end{table*}


Table~\ref{tab:discussion_utbench} presents the performance of LLMs on the internal dataset. For simplicity, we select several representative models from both fine-tuned and instruction-based models that perform well on Defects4J dataset. 
Among fine-tuned models, DeepSeek-Coder-6b achieves the highest performance with a correct rate of 31.14\% and the lowest build error rate of 29.64\%. The encoder-only model UniXcoder performs the lowest, with a correct rate of only 5.39\%, while encoder-decoder models like CodeT5-large and CodeT5p-770m have the lowest syntax error rate at 6.89\% and 3.74\%, respectively. Overall, results on the internal dataset for fine-tuned models closely align with their performance on Defects4J.
In instruction-based models, GPT-3.5 achieves the best results, ranking highest in correct rate, passing rate, and syntax error rate, at 52.54\%, 54.34\%, and 1.20\%, respectively. Llama3.1-70b and Qwen2-72b generally perform well, occupying second and third places across metrics, with Llama3.1-70b achieving a correct rate of 51.35\% and Qwen2-72b at 42.37\%.  

Instruction-based models generally achieve higher correct rates and lower syntax error rates than fine-tuned models, especially when the model size exceeds 70 billion parameters, showing trends similar to Defects4J dataset. 
For clarity, we analyze three model groups: (1) fine-tuned models around 7 billion parameters (e.g., DeepSeek-Coder-6b, CodeLlama-7b, StarCoder2-7b, StarCoderBase-7b, CodeGen-6b), (2) instruction-based models around 7 billion parameters (e.g., Llama3.1-8b, Llama3-8b, Qwen2-7b), and (3) instruction-based models exceeding 70 billion parameters (e.g., GPT-3.5, Llama3.1-70b, Qwen2-72b).
First, the 7-billion-parameter fine-tuned models achieve correct rates between 23.05\% and 31.14\%, averaging 25.27\%. 
Second, the 7-billion-parameter instruction-based models show a modest improvement, with correct rates between 21.71\% and 32.34\% and an average of 27.50\%, marking an 8.82\% increase. 
Third, instruction-based models with over 70 billion parameters achieve correct rates between 42.37\% and 52.54\%, with an average of 48.75\%, reflecting substantial improvements of 92.92\% over the 7-billion fine-tuned models and 77.27\% over the 7-billion instruction-based models.

Overall, these results highlight two main insights. First, among models with similar parameter counts, the performance gap between fine-tuned and instruction-based models is relatively modest. Second, as instruction-based models grow in size and capability, the prompt engineering approach shows strong potential to achieve even higher performance levels.


\summary{
While large-scale instruction-based models outperform fine-tuned models overall, instruction-based models with the same parameters scale show slight gains over fine-tuned counterparts. Instruction-based models show promise as they scale in size and capability.
}
